% Chapter: model V2 - the first-principles results inside the device model, the off-state floors, the held-out
% 6.3 nm test and the ultrathin films. Plain-language style as in ch:dft.
\newenvironment{v2plain}{\par\medskip\noindent\textit{In plain words.}\ }{\par\medskip}

\chapter{Model V2: first-principles inputs, off-state floors and a held-out test}
\label{ch:v2}

\Cref{ch:dft} computed, in one consistent setup, the quantities that the device model had so far borrowed from the
literature: how far the conduction band moves up when the film becomes thin, how the gap opening is shared between the
conduction and the valence band, and how heavy the electrons become. This chapter puts those numbers into the device
model (model V2) and asks four questions:
\begin{description}[leftmargin=1.2em]
\item[V1] What produces the flat off-state currents of the measured transistors, which model V1 does not contain?
\item[V2] How much do the first-principles inputs change the calibrated curves?
\item[V3] Does the model, with nothing adjusted to the 6.3~nm film, predict that film?
\item[V4] What would 1, 0.75 and 0.5~nm films do?
\end{description}

\begin{v2plain}
The device simulator needs a handful of material numbers for every film thickness. Until now the numbers describing
how thinning changes the electron energies came from other people's calculations. We computed them ourselves
(\cref{ch:dft}) and here swap them in. Before running anything we wrote down what each test should show and what
would count as passing, and published that on the project's public repository. That way the results cannot be
bent after the fact.
\end{v2plain}

% =====================================================================================================================
\section{How the work was organised}
\label{sec:v2-method}

The work followed six steps fixed with the user on 5~October~2026:
\begin{itemize}
\item \textbf{T0:} checks without ATLAS.
\item \textbf{T1:} the off-state floor.
\item \textbf{T2:} the first-principles laws in the two calibration films (2.0 and 13.2~nm).
\item \textbf{T3:} the held-out 6.3~nm test.
\item \textbf{T4:} the ultrathin films.
\item \textbf{T5:} this chapter.
\end{itemize}
The user raised the ATLAS launch budget from 57 to 71 for these steps.

\paragraph{Pre-registration.} Every pass/fail criterion, and every predicted value that could be computed without
ATLAS, was written down before the runs. That record is \file{analysis_2026-10-06_V2/T0_PRECHECKS.md}, together with
the SHA-256 hashes of all inputs in \file{T0_REGISTRATION.txt}. It was committed to the public repository as commit
\texttt{7b75ee6} before the first V2 launch.

\paragraph{A correction to the pre-registration.} The first V2 run revealed an error in one of the pre-registered
calculations (\cref{sec:v2-correction}). The correction was written, hashed and published (commit \texttt{a914165})
before the result of the held-out test existed. The registered files were kept unchanged, and both the registered and
the corrected predictions are compared with ATLAS below.

% =====================================================================================================================
\section{The off-state floors}
\label{sec:v2-floor}

\begin{v2plain}
When the measured transistors are switched off, a small current still flows. It does not depend on the gate voltage,
and it is about 1400 times larger in the 13.2~nm film than in the 2~nm film. The simulated transistors switch off
almost completely (about $10^{-19}$~A/\textmu m), so this current is missing from the model. We first asked which
explanations are even possible, then computed whether the remaining one gives the right numbers.
\end{v2plain}

\paragraph{What the data show.} \Cref{tab:v2-floors} lists the floors, taken as the median current between $-2$ and
$-0.5$~V.
\begin{itemize}
\item \textbf{13.2~nm:} a constant 1.9~nA, flat within $\pm3$\,\% over 2.5~V of gate voltage.
\item \textbf{6.3~nm:} flat but noisy ($\pm30$\,\%, mostly random).
\item \textbf{2~nm:} rises smoothly from 0.05 to 1.5~pA. This looks like a charging transient at the start of the
sweep on top of a setup-level floor.
\end{itemize}
From 2 to 13.2~nm the floors grow as $t^{3.85}$.

\begin{table}[htbp]\centering\small
\caption{Measured off-state floors (median over $V_G=-2\ldots-0.5$~V; device width 290~\textmu m) and their shape
(\file{analysis_2026-10-06_V2/t0_floor_data.py}).}\label{tab:v2-floors}
\begin{tabular}{lllp{7.2cm}}\toprule
Film & Floor (A/\textmu m) & Device total & Shape over $-3\ldots-0.5$~V\\\midrule
2.0~nm & $4.61\times10^{-15}$ & 1.3~pA & rises 0.05\,$\to$\,1.5~pA, smooth (start-of-sweep transient on a setup-level floor)\\
6.3~nm & $1.53\times10^{-13}$ & 44~pA & flat, $\pm30$\,\% random noise\\
13.2~nm & $6.60\times10^{-12}$ & 1.9~nA & flat within $\pm3$\,\%\\\bottomrule
\end{tabular}
\end{table}

\paragraph{What is excluded.}
\begin{itemize}
\item \textbf{Gate leakage.} It would be largest at $-3$~V and would change at least threefold over the range. The
6.3 and 13.2~nm floors are flat, and the 2~nm floor is \emph{smallest} at $-3$~V.
\item \textbf{A current range fixed by the on-current.} The 2 and 6.3~nm devices have similar on-currents but floors
33 times apart.
\item \textbf{A conducting layer at the back surface that the gate still controls.} Keeping it flat within 3\,\% over
2.5~V would need about $2\times10^{16}$~cm$^{-2}$\,eV$^{-1}$ of pinning states.
\end{itemize}
What remains is a drain-to-source current outside the reach of the gate, or an instrument effect.

\paragraph{The ungated-strip calculation.} Suppose part of the film lies outside the gate's control. Its sheet
conductance follows from the model's own material values. Charge neutrality fixes the Fermi level $E_F$:
\begin{equation}
N_d\,t \;=\; t\left[n(E_F) + \int g_{\mathrm{tail}}(E) f(E)\,\mathrm{d}E + \int g_{\mathrm{deep}}(E) f(E)\,\mathrm{d}E\right]
 + N_{\mathrm{s}}(E_F) - Q_{\mathrm{s}},
\label{eq:v2-neutral}
\end{equation}
\begin{itemize}
\item $n = N_c F_{1/2}$ is the free-electron density;
\item $g_{\mathrm{tail}}$ and $g_{\mathrm{deep}}$ are the V1 tail and deep acceptor densities of states (\cref{sec:par-eg,sec:par-nc});
\item $N_{\mathrm{s}}$ counts electrons in surface acceptor states, and $Q_{\mathrm{s}}$ is a fixed surface charge.
\end{itemize}
The strip conductance is then $G = q\,n\,\mu_0\,t\,(W/L)_{\mathrm{par}}$. Its one geometric factor $(W/L)_{\mathrm{par}}$
is fitted on the 13.2~nm floor; the 2 and 6.3~nm floors are then predicted. The band bending across the strip from a
surface sheet is a few meV, much less than $k_BT$, so a uniform film is an accurate description.

Three variants were computed (\cref{fig:v2-t0}a):
\begin{itemize}
\item \textbf{N:} no surface states.
\item \textbf{D:} one surface carrying the V1 interface acceptor sheet ($6.4\times10^{10}$~cm$^{-2}$ in total).
\item \textbf{Q:} as D, plus the fixed charge of the channel interface ($1.73\times10^{12}$~cm$^{-2}$).
\end{itemize}

\begin{table}[htbp]\centering\small
\caption{Ungated-strip variants against the measured floors (one geometric factor fitted on 13.2~nm). Criterion set in
the plan: both predicted floors within a factor of 3.}\label{tab:v2-strip}
\begin{tabular}{llccc}\toprule
Variant & Surfaces & Scaling & 2~nm pred/meas & 6.3~nm pred/meas\\\midrule
N & none & $t^{2.99}$ & 5.07 & 1.21\\
D & one surface with the V1 interface acceptors & $t^{4.24}$ & \textbf{0.48} & \textbf{0.77}\\
Q & D plus $Q_f$ of the channel interface & $t^{1.25}$ & 134 & 4.76\\\bottomrule
\end{tabular}
\end{table}

Only variant D passes (\cref{tab:v2-strip}). How firm is that?
\begin{itemize}
\item \textbf{6.3~nm:} the prediction stays within $\times0.37$--$0.99$ when the assumed interface-state density or the
donor density is halved or doubled.
\item \textbf{2~nm:} the prediction swings from $\times0.06$ to $\times1.9$, so the 2~nm agreement is not a strong
test. The 2~nm floor also sits at the setup floor, where an under-prediction would not contradict the data.
\end{itemize}

\paragraph{What V2 adopts.} The total drain current is
\begin{equation}
I_D^{\mathrm{total}}(V_G) = I_D^{\mathrm{ATLAS}}(V_G) + G_{\mathrm{D}}(t)\,V_D ,
\label{eq:v2-strip}
\end{equation}
with $(W/L)_{\mathrm{par}} = 1.03\times10^{-4}$ per \textmu m of width. A gate-independent parallel conductance at
fixed $V_D$ simply adds, so \cref{eq:v2-strip} is exact and costs no launch. Its status is \prov{HYPOTHESIS}. The
location and the interface of the strip are not determined; they need the gate and source currents, a channel-length
series and the layout (\cref{sec:v2-data}).

\begin{figure}[htbp]\centering
\includegraphics[width=\linewidth]{v2_t0_floor_and_laws}
\caption{Step T0 (no ATLAS).
(a) Measured off-state floors (squares) and the three ungated-strip variants, each fitted on 13.2~nm. The band shows
variant D with the interface-state density or the donor density halved or doubled.
(b) Conduction-band shift relative to bulk: V1 law (literature, all of the gap opening in the conduction band) and V2
laws through the DFT points of this work. Dotted lines mark the measured films.
(c) Rigid shift of each calibrated transfer curve when the V1 laws are replaced by the V2 laws, before re-calibration.}
\label{fig:v2-t0}
\end{figure}

% =====================================================================================================================
\section{First-principles inputs in the device model}
\label{sec:v2-dft}

\begin{v2plain}
The DFT results change three numbers in the device model: how far the conduction band rises in a thin film, how much
the gap opens, and how heavy the electrons are. Of these, only the conduction-band rise matters for the transfer curve,
and it changes by less than one tenth of an electronvolt. That moves each calibrated curve sideways by 10 to 75~mV.
\end{v2plain}

\paragraph{New laws.} Fitted through the DFT points of \cref{ch:dft} (PBE):
\begin{align}
\Delta E_c(t) &= 0.78504\,t^{-1.50400}~\text{eV}, &
\Delta E_g(t) &= 0.83997\,t^{-1.34569}~\text{eV}, &
\Delta m^*(t) &= 0.11416\,t^{-1.29453}\,m_0 ,
\label{eq:v2-laws}
\end{align}
with $t$ in nm and $t\geq0.95$~nm. The DFT values behind them are:

\begin{table}[htbp]\centering\small
\caption{DFT values behind \cref{eq:v2-laws} (PBE, relative to bulk).}\label{tab:v2-dftpoints}
\begin{tabular}{lccc}\toprule
Film & $\Delta E_c$ (eV) & $\Delta E_g$ (eV) & $\Delta m^*$ ($m_0$)\\\midrule
0.95~nm & 0.848 & 0.900 & $+0.122$\\
1.98~nm & 0.281 & 0.335 & $+0.047$\\\bottomrule
\end{tabular}
\end{table}

Two points about these values:
\begin{itemize}
\item \textbf{Partition.} The conduction-band share of the gap opening is 0.84 at 2~nm and 0.94 at 1~nm. V1 assumed
1.0, i.e.\ all of the opening in the conduction band.
\item \textbf{Functional.} The HSE06/PBE ratio of the gap opening is 1.06 (range 1.06--1.11), so PBE is kept as the
central value, as the DFT protocol pre-registered.
\end{itemize}
Below 0.95~nm the laws are not extrapolated; the 0.75 and 0.51~nm films enter as tabulated DFT points (\cref{sec:v2-thin}).

\paragraph{Why the effect is a rigid shift.} Every trap level in the decks is referenced to the conduction-band edge,
and the source and drain are ideal Ohmic contacts. A change of the electron affinity therefore moves the whole
transfer curve along $V_G$ by exactly the change of $\Delta E_c$. A change of the front fixed charge is also a rigid
shift, $\Delta V = -q\,\Delta Q_f/C_{ox}$ (verified to 1~mV in V1). The drain current is linear in a uniform mobility.

The mass change enters only through $N_c$. Its effect is taken from the isolation runs 0050/0051 ($m^*\times1.3$
shifts $V_{th,cc}$ by $-0.043$~V at 2~nm) and is below 2~mV here.

So the V2 curves follow from the calibrated V1 runs by exact transformations, before any launch
(\cref{tab:v2-shift}, \cref{fig:v2-t0}c).

\begin{table}[htbp]\centering\small
\caption{Conduction-band shift V1\,$\to$\,V2 and the resulting rigid shift of each calibrated curve.}\label{tab:v2-shift}
\begin{tabular}{lccc}\toprule
Film & $\Delta E_c$ V1 (eV) & $\Delta E_c$ V2 (eV) & Curve shift\\\midrule
2.0~nm & 0.3533 & 0.2768 & $-75$~mV\\
6.3~nm & 0.0724 & 0.0493 & $-24$~mV\\
13.2~nm & 0.0261 & 0.0162 & $-10$~mV\\\bottomrule
\end{tabular}
\end{table}

\paragraph{V2 calibration.} Re-calibration follows the V1 protocol exactly:
\begin{itemize}
\item \textbf{One shared front fixed charge}, chosen by least squares on the 2 and 13.2~nm curves. It moves every curve
by $+62$~mV, giving $Q_f = 1.38345\times10^{12}$~cm$^{-2}$ (V1: $1.73\times10^{12}$).
\item \textbf{One band mobility per calibration film}, scaled exactly to the measured on-current.
\end{itemize}

% =====================================================================================================================
\section{The correction to the pre-registered transformation}
\label{sec:v2-correction}

The exact transformation moves a calibrated ATLAS curve along $V_G$. The registered version clipped the curve at the
end of the simulated sweep (3~V), so a curve moved to the \emph{left} was flat over its last $|\Delta V|$.
\begin{itemize}
\item \textbf{2~nm} (net shift $-13$~mV): the first V2 run gave an on-current 0.83\,\% above the registered prediction,
which revealed the error.
\item \textbf{6.3~nm} (net shift $-0.26$~V): the clipping affected the registered on-state values.
\item \textbf{Not affected:} curves moved to the right (13.2~nm) and the subthreshold region, including $V_{th,cc}$.
\end{itemize}
The fix extrapolates the on-state linearly from the last three simulated points. It was registered and published
before the 6.3~nm result existed (\file{T0_CORRECTION.md}, \file{T0_CORRECTION_REGISTRATION.txt}).

% =====================================================================================================================
\section{ATLAS V2 anchors (T2)}
\label{sec:v2-anchors}
\textit{[filled in from \file{T2_T3_RESULTS.md} after runs 0058--0060]}

\section{The held-out 6.3~nm test (T3)}
\label{sec:v2-test}
\textit{[filled in from \file{T2_T3_RESULTS.md}]}

\section{Ultrathin films (T4)}
\label{sec:v2-thin}
\textit{[filled in after the 0.75 and 0.51~nm DFT films and the scenario runs]}

\section{Data that would decide the open questions}
\label{sec:v2-data}

\begin{v2plain}
Most of the remaining uncertainty cannot be removed by more simulation; it needs a few extra measurements on the
existing wafers. Each item below names the question it would settle.
\end{v2plain}

\paragraph{Off-state floor (\cref{sec:v2-floor}).}
\begin{itemize}
\item \textbf{Gate and source currents recorded with every sweep.} If $I_S \approx -I_D$ and $I_G \approx 0$ in the
off-state, the floor flows from drain to source. If $I_G \approx -I_D$, it is gate leakage.
\item \textbf{A channel-length series at one thickness.} A floor that scales as $1/L$ runs through the channel region;
a floor independent of $L$ runs through a parasitic path elsewhere.
\item \textbf{The instrument settings} (current range, integration time, delay), to separate a real current from the
offset of a range.
\item \textbf{The layout}: whether the film is patterned to the channel and whether the TiN gate covers the whole
substrate.
\end{itemize}

\paragraph{Thickness dependence and the 6.3~nm film (\cref{sec:v2-test}).}
\begin{itemize}
\item \textbf{Several devices per thickness.} The device-to-device spread decides whether a 0.25~V difference is a
model failure or ordinary variation.
\item \textbf{The method and uncertainty of the thickness measurement}, for example X-ray reflectivity.
\item \textbf{Transfer curves at two or more temperatures.} The 65 and 85\,$^\circ$C curves of the 2~nm device
already exist (\citet{Januar2026}, SI Fig.~S12).
\end{itemize}

\paragraph{The first-principles inputs (\cref{sec:v2-dft}).}
\begin{itemize}
\item \textbf{The optical gap against thickness} (UV--Vis or ellipsometry): the measured $\Delta E_g(t)$.
\item \textbf{X-ray photoelectron valence-band spectra}: the valence-band shift, i.e.\ the computed 0.84--0.94 share.
\item \textbf{Hall measurements on unpatterned films}: carrier density and band mobility, which V2 still fits per film.
\end{itemize}
